\section{Preliminaries}
\label{sec:preliminaries}

This section fixes notation and recalls the formal syntax of the lattice primitives, the BB-tran rational hash,
non-interactive zero-knowledge arguments, blind signatures, anonymous rate-limited credentials, and
pseudorandom functions.  We also record the imported non-blind signature-security statements from
Dubois--Kloo\ss--Lai--Woo (DKLW) that are used later for the hash-and-preimage signature layer.

\subsection{Notation}
Let $\lambda$ denote the security parameter and let $\negl(\lambda)$ be any positive function negligible in
$\lambda$.  We write PPT for probabilistic polynomial time.  For integers $a<b$, let
$[a,b]=\{a,a+1,\ldots,b\}$ and use $[b]$ as shorthand for $[1,b]$.  Vectors are written in bold lower-case
letters and matrices in bold upper-case letters.

We work over the cyclotomic ring
\[
  \R=\Z[X]/(X^N+1)
\]
for power-of-two $N$, and its quotient
\[
  \Rq=\R/q\R.
\]
We identify $\Rq$ with polynomials of degree $<N$ with coefficients in $\Z_q$.  For $a\in\Z$, we write
$\langle a\rangle_q\in\{0,\ldots,q-1\}$ for reduction modulo $q$.  Whenever a field element is interpreted as a
signed integer, we use its centered representative in $[-q/2,q/2]$.  For $x\in\Rq$, the notation
$\|x\|_\infty\le B$ means that every coefficient of $x$ lies in $[-B,B]$; the same convention is used
component-wise for vectors over $\Rq$.

The full signed message is written
\[
  \vecmu=\vecm\Vert \veck,
\]
where $\vecm$ encodes holder attributes and $\veck$ is the secret PRF key used for rate limiting.  The key
is sampled from a finite key space $\mathcal K_{\mathsf{key}}\subseteq \Fq^{\ell_{\mathsf{key}}}$ fixed by the
public parameters.  In the bounded-key instantiation used by the present relations, we require
\[
  \mathcal K_{\mathsf{key}}\subseteq
  \{\veck\in\Fq^{\ell_{\mathsf{key}}}:\|\veck\|_\infty\le \BoundMsg\},
\]
and the PRF assumption is made for a uniform key from this same key space.  The public coefficient bound
for holder attributes and protocol randomness is denoted $\BoundMsg$, and the short-preimage bound for
signatures is denoted $\BoundSig$.

Whenever the CRT decomposition $\Rq\cong\prod_i F_i$ is used, we write
\[
  (x_i)_i\Had(y_i)_i := (x_i y_i)_i,
  \qquad
  (x_i)_i\Hdiv (y_i)_i := (x_i y_i^{-1})_i
\]
whenever every $y_i\neq 0$.  Thus $\Had$ and $\Hdiv$ denote slotwise multiplication and slotwise division in
CRT coordinates.

\subsection{Centering map and shared randomness}
\label{subsec:centering}
The issuance protocol combines holder-side and issuer-side bounded randomness by a coefficient-wise
centering operation.

\begin{definition}[Centering map]\label{def:centering}
Fix the coefficient bound $\BoundMsg$ and let
\[
  M_{\mathrm{ctr}}:=2\BoundMsg+1.
\]
For an integer $a\in\Z$, define
\[
  \centerfunc(a)\in[-\BoundMsg,\BoundMsg]
\]
to be the unique representative of $a$ modulo $M_{\mathrm{ctr}}$ in that interval.  Extend $\centerfunc$ coefficient-wise
from integers to ring elements of $\Rq$ and then component-wise to vectors over $\Rq$.
\end{definition}

\begin{lemma}[Uniformity of centered sums]\label{lem:center-uniform}
Let $x^H\in[-\BoundMsg,\BoundMsg]$ be fixed and let $x^I$ be uniform over $[-\BoundMsg,\BoundMsg]$.
Then
\[
  x=\centerfunc(x^H+x^I)
\]
is uniform over $[-\BoundMsg,\BoundMsg]$.  The same holds coefficient-wise for ring elements and vectors over
$\Rq$.
\end{lemma}

\begin{proof}
The map $x^I\mapsto \centerfunc(x^H+x^I)$ is a permutation of the interval $[-\BoundMsg,\BoundMsg]$ because
it is translation modulo $M_{\mathrm{ctr}}$.  Therefore it preserves the uniform distribution.  Applying this
argument coefficient-wise proves the ring and vector statements.
\end{proof}

\subsection{Commitment schemes and Ajtai commitments}
\label{subsec:prelim-commit}
\label{subsec:commitments}

\begin{definition}[Commitment scheme]\label{def:commitment-scheme}
A non-interactive commitment scheme is a tuple of PPT algorithms
\[
  \ComScheme=(\ComSetup,\Commit,\ComVerify)
\]
with the following syntax.
\begin{itemize}[leftmargin=1.25em]
  \item $\ComPublicParam\gets \ComSetup(\SecParam,\PublicParamsGen)$ outputs public commitment parameters.
  \item $(c,d)\gets \Commit(m)$ takes a message $m$ and outputs a commitment $c$ and an opening $d$.
  \item $\ComVerify(c,m,d)\in\{0,1\}$ deterministically checks whether $d$ is a valid opening of $c$ to $m$.
\end{itemize}
The scheme is correct if honestly generated commitments always verify.  It is computationally binding if no PPT
adversary can open one commitment to two distinct messages except with negligible probability.  It is hiding if
commitments to any two messages are computationally indistinguishable.
\end{definition}

\subsubsection{Ajtai commitment used in ARC-SPRUCE.}
Fix message length $\ell_\mu$, holder-randomness lengths $\ell_{r_0},\ell_{r_1}$, commitment-only randomness length
$\ell_{\bar r}$, and output length $\ell_{\ComScheme}$.  Let
\[
  \ell_{\mathrm{in}}:=\ell_\mu+\ell_{r_0}+\ell_{r_1}+\ell_{\bar r}.
\]
The public commitment matrix is
\[
  \mACom\getsunif \Rq^{\ell_{\ComScheme}\times \ell_{\mathrm{in}}}.
\]
When useful, we block-parse it as
\[
  \mACom=[\mA_\mu\mid \mA_{r_0}\mid \mA_{r_1}\mid \mA_{\bar r}],
\]
so that the message part and the three randomness blocks are explicit.

A holder committing to the full message $\vecmu$ and to holder-side rational-hash randomness
$(r_0^H,r_1^H)$ samples additional commitment-only randomness $\bar r$ and computes
\begin{equation}
  \vecc
  =
  \mACom[\vecmu\Vert r_0^H\Vert r_1^H\Vert \bar r]^\top.
  \label{eq:ajtai-commit-current}
\end{equation}
Equivalently,
\[
  \vecc=\mA_\mu\vecmu+\mA_{r_0}r_0^H+\mA_{r_1}r_1^H+\mA_{\bar r}\bar r.
\]
The opening is $(\vecmu,r_0^H,r_1^H,\bar r)$.  Verification checks the linear equation
\eqref{eq:ajtai-commit-current} together with the coefficient bounds on all opened values.

\begin{definition}[Module-SIS and Module-ISIS over $\Rq$]\label{def:sis-isis}
Fix integers $n,m\ge 1$, modulus $q\ge 2$, and bound $\beta>0$.
\begin{itemize}[leftmargin=1.25em]
  \item \textbf{$\mathsf{M\text{-}SIS}_\infty(n,m,q,\beta)$.} Given $\mA\in\Rq^{n\times m}$, find a non-zero
  $\mathbf{s}\in \R^m$ such that $\mA\mathbf{s}=0$ in $\Rq^n$ and $\|\mathbf{s}\|_\infty\le \beta$.
  \item \textbf{$\mathsf{M\text{-}ISIS}_\infty(n,m,q,\beta)$.} Given $(\mA,\vect)\in\Rq^{n\times m}\times\Rq^n$,
  find $\mathbf{s}\in \R^m$ such that $\mA\mathbf{s}=\vect$ in $\Rq^n$ and
  $\|\mathbf{s}\|_\infty\le \beta$.
\end{itemize}
\end{definition}

\begin{lemma}[Binding of the Ajtai commitment]\label{lem:ajtai-binding}
If a PPT adversary opens one Ajtai commitment to two distinct bounded vectors with non-negligible
probability, then there exists a PPT adversary that solves the corresponding
$\mathsf{M\text{-}SIS}_\infty$ instance with essentially the same probability and runtime.
\end{lemma}

\begin{proof}
Suppose two distinct valid openings
\[
  (\vecmu,r_0^H,r_1^H,\bar r)\neq (\vecmu',r_0^{H\prime},r_1^{H\prime},\bar r')
\]
verify the same commitment $\vecc$.  Subtracting the two opening equations gives
\[
  \mACom\bigl([
    \vecmu-\vecmu'\Vert
    r_0^H-r_0^{H\prime}\Vert
    r_1^H-r_1^{H\prime}\Vert
    \bar r-\bar r']^\top
  \bigr)=0.
\]
The difference vector is non-zero and coefficient-wise bounded by $2\BoundMsg$, hence it is an
$\mathsf{M\text{-}SIS}_\infty$ solution.
\end{proof}

\begin{lemma}[Statistical hiding of the Ajtai commitment]\label{lem:ajtai-hiding}
Assume $q$ is unramified in $\R$ and write the factorization
\[
  X^N+1 \equiv \prod_{j=1}^{f_{\mathrm{fac}}}\Phi_j(X) \pmod q,
  \qquad
  \deg \Phi_j=d_{\mathrm{fac}},
  \qquad
  N=f_{\mathrm{fac}}d_{\mathrm{fac}}.
\]
Let the commitment-only randomness $\bar r$ be sampled coefficient-wise from $[-B_r,B_r]$ with
$B_r<q/2$.  Define
\[
  \varepsilon_{\mathrm{Com}}
  :=
  \frac12
  \sqrt{
    \left(
      \left(
        \frac{q^{\ell_{\ComScheme}}}{(2B_r+1)^{\ell_{\bar r}}}
      \right)^{d_{\mathrm{fac}}}
      +1
    \right)^{f_{\mathrm{fac}}}
    -1
  }.
\]
Then the commitment distribution is statistically $\varepsilon_{\mathrm{Com}}$-close to a distribution
independent of the committed vector $(\vecmu,r_0^H,r_1^H)$.  The commitment parameters are
$\lambda$-hiding when $\varepsilon_{\mathrm{Com}}\le 2^{-\lambda}$.  In the bounded instantiation used here,
one may take $B_r=\BoundMsg$.
\end{lemma}

\begin{proof}
For fixed $(\vecmu,r_0^H,r_1^H)$, the commitment is a public offset plus the linear image
$\mA_{\bar r}\bar r$.  Applying the bounded-uniform cyclotomic leftover-hash lemma with input dimension
$\ell_{\bar r}$, output dimension $\ell_{\ComScheme}$, and coefficient radius $B_r$ gives the displayed
statistical-distance bound.  Adding the fixed offset does not change statistical distance.
\end{proof}

\subsection{Lattice trapdoors and rational hashing}
\label{subsec:prelim-vsis}
We recall the interface used by the hash-and-preimage signature layer.  Additional lattice-assumption
syntax, such as vSIS, hinted-vSIS, and GenISIS$_f$, is deferred to the appendix.

\begin{definition}[Admissible lattice trapdoor suite]\label{def:trapdoor-suite}
A lattice trapdoor suite over $\Rq$ consists of PPT algorithms
\[
  (\Setup,\TrapGen,\SampPre)
\]
together with a public bound $\BoundSig$.  We say that the suite is admissible for the parameters used in
this paper if the following interface holds.
\begin{itemize}[leftmargin=1.25em]
  \item $\PublicParamsTrapdoor\gets\Setup(\SecParam,\PublicParamsGen)$ outputs dimensions $(n,m)$, modulus
  $q$, sampler parameters, and the accepted shortness bound $\BoundSig$.
  \item $(\mA,\trapdoor)\gets\TrapGen(\PublicParamsTrapdoor)$ outputs a public matrix
  $\mA\in\Rq^{n\times m}$ together with a trapdoor $\trapdoor$, and the distribution of $\mA$ is
  statistically or computationally close to uniform over $\Rq^{n\times m}$ as required by the invoked
  signature theorem.
  \item On input $(\PublicParamsTrapdoor,\trapdoor,T)$ with $T\in\Rq^n$, the sampler
  $\SampPre(\trapdoor,T)$ outputs $\vecu\in\R^m$ such that $\mA\vecu=T$ in $\Rq^n$, except with
  negligible failure probability, and $\|\vecu\|_\infty\le \BoundSig$ except with negligible probability.
  \item The distributional hiding property required by the DKLW hash-and-preimage signature theorem holds for
  the produced preimages at the chosen Gaussian width and public bound.
\end{itemize}
This definition is an interface assumption for the signature-layer import.  Appendix~\ref{app:gaussian-sampler}
describes the NTRU/Antrag-based candidate used for parameter tuning, but Antrag alone is not cited here as a
proof of this full admissible-suite interface.
\end{definition}

\begin{definition}[Trapdoor rational hash]\label{def:trapdoor-rational-hash}
A trapdoor rational hash scheme over $\Rq$ is a tuple
\[
  \TDRH=(\Setup,\TrapGen,\Eval,\SampPre)
\]
consisting of the lattice trapdoor suite above together with a public evaluation algorithm.  We extend
$\Setup$ so that $\PublicParamsTrapdoor$ also contains a public hash description $B$, a randomness space
$\mathcal X$, a message space $\mathcal M$, and an admissibility error bound $\delta$.  For fixed $B$, the
algorithm
\[
  \Eval(\PublicParamsTrapdoor,\mu,\chi)
\]
takes a message $\mu\in\mathcal M$ and witness randomness $\chi\in\mathcal X$ and outputs either a target
$T\in\Rq^n$ or $\bot$ if the input is inadmissible.  For a public description $B$ contained in
$\PublicParamsTrapdoor$, define the admissible message subspace
\[
  \mathcal M_B:=\left\{\mu\in\mathcal M:
    \Pr_{\chi\gets \mathcal X}\bigl[\Eval(\PublicParamsTrapdoor,\mu,\chi)=\bot\bigr]\le \delta
  \right\}.
\]
Correctness requires that for every $\mu\in\mathcal M_B$ and every $\chi\in\mathcal X$ such that
$\Eval(\PublicParamsTrapdoor,\mu,\chi)=T\neq\bot$, the output
$\vecu\gets\SampPre(\trapdoor,T)$ satisfies
\[
  \mA\vecu=T
  \qquad\text{and}\qquad
  \|\vecu\|_\infty\le \BoundSig
\]
except with negligible probability.
\end{definition}

The generic trapdoor-rational-hash interface does not by itself imply that the public target hides the signed
message.  In the concrete BB-tran instantiation used in this paper, target hiding is obtained later from a
cyclotomic leftover-hash lemma applied to the long linear-side randomizer.

\paragraph{The BB-tran instantiation.}
We instantiate the rational hash with the BB-tran family from DKLW~\cite[Table~1]{EPRINT:DKLW25}.  In
that work the generic expression is
\[
  h_{\mu,\chi}(\tilde b^\top)
  =
  (u^\top,x_0^\top)\tilde b_0+\frac{1}{\tilde b_1-x_1}.
\]
The translation used in ARC-SPRUCE is
\[
  u\leftrightarrow \vecmu,\qquad
  x_0\leftrightarrow r_0,\qquad
  x_1\leftrightarrow r_1.
\]
To account for the affine term used in the protocol, we view the DKLW linear side as having an implicit
leading coordinate:
\[
  (1,\vecmu^\top,r_0^\top)\widehat B_0
  =
  B_0+B_1\vecmu+B_2r_0.
\]
Thus the public description is
\[
  B=(B_0,B_1,B_2,B_3),
\]
where $B_0,B_3\in\Rq^n$ and $B_1,B_2$ are public linear maps into $\Rq^n$.  The inverse-side
randomizer $r_1\in\Rq$ is broadcast to all target coordinates; for readability we write $B_3-r_1$ for
$B_3-\mathbf 1_n r_1$.  If an implementation instead uses a vector-valued $r_1\in\Rq^n$, that is a
vectorized variant and its DKLW predicate and security side conditions must be assumed separately.

For message $\vecmu$ and randomness $(r_0,r_1)$, define the linear side
\begin{equation}
  L_B(\vecmu,r_0):=B_0+B_1\vecmu+B_2r_0.
  \label{eq:htran-linear-side}
\end{equation}
The input $(B,\vecmu,r_0,r_1)$ is admissible if $B_3-\mathbf 1_n r_1\in\Rq^{\times n}$, equivalently if
every CRT slot of every coordinate of $B_3-\mathbf 1_n r_1$ is non-zero.  When admissible, let
\[
  Z=(B_3-\mathbf 1_n r_1)^{-1}
\]
with inversion understood coordinate-wise in CRT form.  The BB-tran target is then
\begin{equation}
  T=h_{\mathrm{tran},\vecmu,(r_0,r_1)}(B):=L_B(\vecmu,r_0)+Z.
  \label{eq:htran-target}
\end{equation}
Equivalently, the witness form is
\begin{equation}
  (B_3-\mathbf 1_n r_1)\Had Z=1,
  \qquad
  T=B_0+B_1\vecmu+B_2r_0+Z,
  \label{eq:htran-witness-form}
\end{equation}
and the eliminated form is
\begin{equation}
  (B_3-\mathbf 1_n r_1)\Had\bigl(T-B_0-B_1\vecmu-B_2r_0\bigr)=1.
  \label{eq:htran-eliminated}
\end{equation}
The witness form is preferred in the protocol statements because it isolates the bilinear inverse check from
the linear identity for the public target.

\begin{lemma}[Eliminated form under admissibility]\label{lem:hash-cleared-admissible}
Let $\Rq\cong\prod_i F_i$ be the CRT decomposition and write $B_3-\mathbf 1_n r_1=(d_i)_i$ and
$T-L_B(\vecmu,r_0)=(z_i)_i$.  If $d_i\neq 0$ for all CRT slots, then
\[
  T=h_{\mathrm{tran},\vecmu,(r_0,r_1)}(B)
  \quad\Longleftrightarrow\quad
  (B_3-\mathbf 1_n r_1)\Had\bigl(T-B_0-B_1\vecmu-B_2r_0\bigr)=1.
\]
\end{lemma}

\begin{proof}
In each CRT slot, the definition is the scalar statement $z_i=d_i^{-1}$, which is equivalent to $d_i z_i=1$
when $d_i\neq 0$.  Reassembling the slots proves the equivalence.
\end{proof}

\begin{definition}[Long target-randomizer]\label{def:htran-long-r0}
Fix an integer $\ell_{r_0}\ge 1$.  The linear-side BB-tran randomizer is sampled as
\[
  r_0\getsunif [-\BoundMsg,\BoundMsg]^{N\ell_{r_0}}\subseteq \Rq^{\ell_{r_0}},
\]
meaning that each coefficient of each ring coordinate of $r_0$ is sampled independently and uniformly from
$[-\BoundMsg,\BoundMsg]$.
\end{definition}

\begin{lemma}[Cyclotomic leftover-hash specialization]\label{lem:htran-lhl}
Assume $q$ is unramified in $\R$ and write
\[
  X^N+1 \equiv \prod_{j=1}^{f_{\mathrm{fac}}}\Phi_j(X) \pmod q,
  \qquad
  \deg \Phi_j=d_{\mathrm{fac}},
  \qquad
  N=f_{\mathrm{fac}}d_{\mathrm{fac}}.
\]
Let $n\ge 1$, let $B_2\getsunif \Rq^{n\times \ell_{r_0}}$, and let $r_0$ be distributed as in
Definition~\ref{def:htran-long-r0}.  Define
\[
  \varepsilon_{r_0}
  :=
  \frac12
  \sqrt{
    \left(
      \left(
        \frac{q^n}{(2\BoundMsg+1)^{\ell_{r_0}}}
      \right)^{d_{\mathrm{fac}}}
      +1
    \right)^{f_{\mathrm{fac}}}
    -1
  }.
\]
Then
\[
  \Delta\bigl((B_2,B_2r_0),(B_2,U)\bigr)\le \varepsilon_{r_0},
  \qquad
  U\getsunif \Rq^n.
\]
In particular, the parameters give $\lambda$-bit target hiding from this step when
$\varepsilon_{r_0}\le 2^{-\lambda}$.
\end{lemma}

\begin{proof}
Apply the bounded-uniform cyclotomic leftover-hash lemma to the linear map
$r_0\mapsto B_2r_0$ with input dimension $\ell_{r_0}$, output dimension $n$, and coefficient radius
$\BoundMsg$.  The displayed expression is the resulting statistical-distance bound for the stated
factorization of $R_q$.
\end{proof}

\begin{theorem}[LHL-based target hiding of $h_{\mathrm{tran}}$]\label{thm:htran-target-hiding}
Assume the hypotheses of Lemma~\ref{lem:htran-lhl}, so the public parameters are honestly generated and
in particular $B_2$ is sampled uniformly at setup.  Let $r_1$ be any fixed admissible inverse-side randomizer
and define, for a bounded message $\mu\in\mathcal M$,
\[
  T_\mu := B_0+B_1\mu+B_2r_0+(B_3-\mathbf 1_n r_1)^{-1}.
\]
Then $T_\mu$ is statistically $\varepsilon_{r_0}$-close to uniform over $\Rq^n$.  Consequently, for any two
bounded
messages $\mu_0,\mu_1\in\mathcal M$,
\[
  \Delta(T_{\mu_0},T_{\mu_1})\le 2\varepsilon_{r_0}.
\]
If $\varepsilon_{r_0}\le 2^{-\lambda}$, this gives
$\Delta(T_{\mu_0},T_{\mu_1})\le 2^{1-\lambda}$.
\end{theorem}

\begin{proof}
By Lemma~\ref{lem:htran-lhl}, the distribution of $B_2r_0$ is statistically $\varepsilon_{r_0}$-close to a
uniform
$U\getsunif\Rq^n$.  For fixed admissible $r_1$ and fixed message $\mu$, the remaining term
\[
  C_{\mu,r_1}:=B_0+B_1\mu+(B_3-\mathbf 1_n r_1)^{-1}
\]
is a deterministic offset, hence
\[
  T_\mu=B_2r_0+C_{\mu,r_1}
\]
is also statistically $\varepsilon_{r_0}$-close to uniform.  Applying the same argument to $\mu_0$ and
$\mu_1$ and using the triangle inequality yields the displayed bound.  No coefficient-norm bound on
$(B_3-\mathbf 1_n r_1)^{-1}$ is used in this hiding argument: the inverse-side term enters only as an
additive offset.
\end{proof}

\begin{corollary}[Protocol distribution of the target randomizer]\label{cor:htran-protocol-r0}
Assume $r_0^H,r_0^I\in\Rq^{\ell_{r_0}}$ are sampled coefficient-wise in $[-\BoundMsg,\BoundMsg]$, with
$r_0^I$ uniform.  Then
\[
  r_0=\centerfunc(r_0^H+r_0^I)
\]
is coefficient-wise uniform over $[-\BoundMsg,\BoundMsg]$ by Lemma~\ref{lem:center-uniform}.  Hence the
protocol distribution of $r_0$ satisfies Definition~\ref{def:htran-long-r0}.  If issuance restarts until
$B_3-\mathbf 1_n r_1$ is admissible, the target-hiding argument remains valid after conditioning on
admissibility because Theorem~\ref{thm:htran-target-hiding} is pointwise in every fixed admissible $r_1$.
\end{corollary}

\subsection{Signatures and imported non-blind security}
\label{subsec:prelim-signatures}
The blind-issuance layer in this paper builds on the non-blind vSIS hash-and-preimage signature layer from
DKLW.  We therefore record the standard signature syntax and the strong unforgeability notions used by that
import.

\begin{definition}[Signature scheme]\label{def:signature-scheme}
A signature scheme for a message space $\mathcal M$ is a tuple of PPT algorithms
\[
  \Sigma=(\Setup,\KeyGen,\Sign,\SigVerify).
\]
It is correct if honestly generated signatures verify except with negligible probability.
\end{definition}

\begin{definition}[Strong existential unforgeability]\label{def:seuf-defs}
Let $\Sigma$ be a signature scheme.  We use the standard notions of strong existential unforgeability under
chosen-message attack (sEUF-CMA), selective-message attack (sEUF-SMA), and strong selective unforgeability
under selective-message attack (sSUF-SMA), exactly as in the DKLW formalization.
\end{definition}

\begin{definition}[vSIS-based hash-and-preimage signature]\label{def:vsis-template}
Let $\TDRH$ be a trapdoor rational hash scheme with parameters
$(B,\mathcal M,\mathcal X,\delta,\BoundSig)$.  The associated vSIS signature scheme
\[
  \vsisscheme=(\KeyGen,\Sign,\SigVerify)
\]
is defined as follows.
\begin{itemize}[leftmargin=1.25em]
  \item $\KeyGen$ runs $(\mA,\trapdoor)\gets\TrapGen(\PublicParamsTrapdoor)$ and outputs
  $(\vk,\sk)=(\mA,\trapdoor)$.
  \item $\Sign(\sk,\mu)$ samples $\chi\gets\mathcal X$ until $T=\Eval(\PublicParamsTrapdoor,\mu,\chi)\neq \bot$,
  then samples $\vecu\gets\SampPre(\trapdoor,T)$ and outputs
  \[
    \sigma=(\chi,\vecu).
  \]
  \item $\SigVerify(\vk,\mu,\sigma)$ for $\sigma=(\chi,\vecu)$ checks that $\chi\in\mathcal X$,
  $\mu\in\mathcal M_B$, $\|\vecu\|_\infty\le \BoundSig$, and
  \[
    \mA\vecu=\Eval(\PublicParamsTrapdoor,\mu,\chi).
  \]
\end{itemize}
\end{definition}

\begin{theorem}[Specialized DKLW security import]\label{thm:dkwl-seuf-sma}
Let $\Sigma_{\mathrm{tran}}$ be the vSIS signature scheme obtained from Definition~\ref{def:vsis-template} by
instantiating the rational family with the BB-tran form of \eqref{eq:htran-target}.  Assume the admissible
trapdoor-suite interface of Definition~\ref{def:trapdoor-suite}, the DKLW predicate and denominator-bound
conditions, the required strong-linear-independence condition for the instantiated rational functions, the
required norm-compatibility conversion between DKLW's norm convention and the public bound
$\BoundSig$, negligible inadmissibility errors, and the corresponding strong hinted-vSIS hypothesis.  Then
$\Sigma_{\mathrm{tran}}$ is strongly existentially unforgeable under selective-message attack.
\end{theorem}

\begin{proof}
This is the BB-tran specialization of the DKLW selective-query security theorem for their generic
vSIS-based signature family.  The statement here is conditional on the DKLW side conditions being
satisfied by the concrete ARC-SPRUCE parameters.
\end{proof}

\begin{corollary}[Conditional imported route to sEUF-CMA]\label{cor:dkwl-seuf-cma}
Suppose, in addition, that the exact hypotheses of the DKLW $\GenISIS_f/\IntGenISIS_f$ lifting theorem
for BB-tran are met, including the required power-of-two cyclotomic setting, the modulus and message-bound
condition $q/2>\gamma$, the dimension condition $m=n\log q+\omega(\log\lambda)$, the smoothing lower
bound on the Gaussian parameter, and the theorem's stated bound loss.  Then the same non-blind
signature layer $\Sigma_{\mathrm{tran}}$ obtains the corresponding strong existential unforgeability guarantee
under adaptive chosen-message attack.
\end{corollary}

\begin{proof}
This is a conditional import of DKLW's GenISIS/IntGenISIS lifting theorem.  It concerns only the non-blind
hash-and-preimage signature layer and does not imply one-more unforgeability of the interactive issuance
protocol.
\end{proof}

\begin{remark}[Scope of the import]\label{rem:dkwl-scope}
Theorems~\ref{thm:dkwl-seuf-sma} and~\ref{cor:dkwl-seuf-cma} concern the non-blind hash-and-preimage
signature layer only.  They do not by themselves prove one-more unforgeability of the blind issuance
protocol.
\end{remark}

\subsection{Non-interactive zero-knowledge arguments of knowledge}
\label{subsec:prelim-nizk}
\label{subsec:nizk-ark}
We use the NIZK-ARK syntax of the SmallWood framework instantiated in the random-oracle model.

\begin{definition}[NIZKs in the random-oracle model]\label{def:nizk-ark}
Let $\mathcal R\subseteq \mathcal X\times \mathcal W$ be a family of relations.  A transparent non-interactive
zero-knowledge argument of knowledge for $\mathcal R$ in the random-oracle model, with oracle $\RO$, is a
pair of algorithms $(\ZKProve^{\RO},\ZKVerify^{\RO})$:
\begin{itemize}[leftmargin=1.25em]
  \item $\ZKProve^{\RO}(w\mid \mathcal R,x)\to \pi$ is probabilistic and, on input a statement-witness pair
  $(x,w)\in\mathcal R$, outputs a proof $\pi$.
  \item $\ZKVerify^{\RO}(\pi\mid \mathcal R,x)\to\{0,1\}$ is deterministic.
\end{itemize}
It satisfies the following properties.

\paragraph{Completeness.}
For every $(x,w)\in\mathcal R$,
\[
  \Pr\bigl[\ZKVerify^{\RO}(\pi\mid \mathcal R,x)=1 : \pi\gets\ZKProve^{\RO}(w\mid \mathcal R,x)\bigr]=1.
\]

\paragraph{Straightline-extractable knowledge soundness.}
There exists a PPT extractor $\Ext$ such that for every PPT adversary $\Adv$,
\[
  \Pr\Bigl[
    \ZKVerify^{\RO}(\pi\mid \mathcal R,x)=1\ \wedge\ (x,w)\notin \mathcal R
  \Bigr]\le \varepsilon,
\]
where $(\pi,Q)\leftarrow \Adv^{\RO}(x)$, $Q$ is the set of random-oracle query/answer pairs made by $\Adv$,
and $w\leftarrow \Ext(Q,\pi)$.

\paragraph{Zero knowledge.}
There exists a PPT simulator $\Sim$, which may program the random oracle, such that for every
$(x,w)\in\mathcal R$ and every PPT distinguisher $\Adv$,
\[
  \Bigl|\Pr[\Adv^{\RO}(\pi_{\mathrm{real}})=1]-\Pr[\Adv^{\RO}(\pi_{\mathrm{sim}})=1]\Bigr|\le \xi,
\]
where
\[
  \pi_{\mathrm{real}}\gets \ZKProve^{\RO}(w\mid \mathcal R,x),
  \qquad
  \pi_{\mathrm{sim}}\gets \Sim^{\RO}(x).
\]
\end{definition}

We instantiate these proofs through the SmallWood PACS/PIOP framework compiled with Fiat--Shamir in the
random-oracle model.

\subsection{Blind signatures}
\label{subsec:blind-sig-defs}
We follow the standard blind-signature syntax.

\begin{definition}[Blind-signature syntax]\label{def:bs-syntax}
A blind signature scheme is a tuple
\[
  \BlindSignature=(\Setup,\IKeyGen,\Issue,\Verify)
\]
where $\Setup$ generates public parameters, $\IKeyGen$ generates issuer keys $(\vk,\sk)$, $\Issue$ is an
interactive protocol between issuer and holder, and $\Verify(\vk,\mu,\sig)$ checks signatures.
\end{definition}

\begin{definition}[Blind-signature correctness]\label{def:bs-correctness}
A blind signature scheme is correct if there exists $\negl$ such that for every message $\mu$,
\[
\Pr\!\left[
\begin{array}{l}
  \PublicParamsBS \gets \Setup(\SecParam);\\
  (\vk,\sk)\gets \IKeyGen(\PublicParamsBS);\\
  \sig\gets \Holder^{\Issue:\langle \Issuer(\sk):\cdot\rangle}(\mu,\vk);\\
  \Verify(\vk,\mu,\sig)=1
\end{array}
\right]
\ge 1-\negl(\lambda).
\]
\end{definition}

\begin{definition}[Blindness]\label{def:bs-blindness}
A blind signature scheme is blind if, for any PPT issuer-side adversary interacting with two honest holders on
challenge messages $\mu_0,\mu_1$, the completed signatures are computationally indistinguishable up to their
output order.
\end{definition}

\begin{definition}[One-more unforgeability]\label{def:bs-one-more}
A blind signature scheme is one-more unforgeable if any PPT adversary participating in at most $Q$ issuance
interactions with the honest issuer cannot output $Q+1$ valid signatures on distinct messages, except with
negligible probability.
\end{definition}

\begin{remark}[Scope of the issuance layer]\label{rem:bs-scope}
The protocol layer studied in Section~\ref{sec:blind-signature} is used as an issuer-view-hiding blind-issuance
protocol together with a raw post-issuance verification relation.  The current draft does not claim that its
completed output instantiates Definition~\ref{def:bs-blindness}, and one-more unforgeability of that
interactive issuance layer remains a separate reduction goal.
\end{remark}

\subsection{Anonymous Rate-Limiting Credentials}
\label{subsec:arc-security}
Rate limiting is part of the ARC syntax: the holder presents a credential together with a public nonce and a
PRF-derived public tag, while the verifier enforces the policy by maintaining state on previously accepted
$(\nonce,\prftag)$ pairs.

\begin{definition}[ARC syntax]\label{def:arc-syntax}
Let $\nonceSpace$ be a finite nonce set.  An anonymous rate-limiting credential scheme is a tuple
\[
  \ARC=(\Setup,\IKeyGen,\Issue,\Show,\Verify)
\]
where $\Issue$ outputs a credential, $\Show(\credential,\nonce)$ outputs a presentation, and
$\Verify(\vk,\presentation,\algstate)$ checks the presentation and updates verifier state.
\end{definition}

\begin{definition}[ARC soundness]\label{def:arc-soundness}
Let $\ARC$ have nonce space $\nonceSpace$.  After at most $Q$ issuance interactions, no PPT adversary should
be able to produce more than $Q|\nonceSpace|$ accepted presentations, except with negligible probability.
\end{definition}

\begin{definition}[Presentation unlinkability]\label{def:pres-unlink}
An ARC scheme has presentation unlinkability if, for distinct fresh nonces, presentations produced from the
same credential are computationally indistinguishable from presentations produced from independently issued
credentials, up to the intended linkage caused by nonce reuse.
\end{definition}

\begin{remark}[Scope of the ARC theorems]\label{rem:arc-scope}
Section~\ref{sec:arc-construction} does not prove Definition~\ref{def:arc-soundness} in full.  It records a
conditional policy-soundness statement and a conditional presentation-unlinkability statement for the ARC
construction induced by the current issuance layer.
\end{remark}

\subsection{Pseudorandom functions}
\label{subsec:prelim-prf}
\label{subsec:prf}

\begin{definition}[Pseudorandom function]\label{def:prf}
Let $\mathcal K,\mathcal D,\mathcal T$ be finite sets.  A family
\[
  \PRF:\mathcal K\times \mathcal D\to \mathcal T
\]
is pseudorandom if oracle access to $\PRF(k,\cdot)$ for uniform $k\getsunif\mathcal K$ is computationally
indistinguishable from oracle access to a uniform random function from $\mathcal D$ to $\mathcal T$.
\end{definition}

We instantiate the PRF with a Poseidon2-style permutation, but the PRF security of the keyed
feed-forward/truncated family below is an explicit assumption rather than a theorem derived from the
Poseidon2 permutation paper.

\begin{definition}[Poseidon2-style permutation]\label{def:poseidon-perm}
Let $q$ be prime, let $d\ge 3$ satisfy $\gcd(d,q-1)=1$, and let $R_F,R_P\in\mathbb N$ with $R_F$ even.  A
Poseidon2-style permutation alternates full external rounds and partial internal rounds, together with public
round constants and invertible linear mixing maps over $\Fq$.
\end{definition}

\begin{definition}[Poseidon2-style PRF]\label{def:poseidon-prf}
For key $\veck\in\mathcal K_{\mathsf{key}}\subseteq\Fq^{\ell_{\mathsf{key}}}$ and public input
$\vecn\in\Fq^{\ell_{\mathsf{nonce}}}$, define
\[
  \PRF(\veck,\vecn)=\Tr_{\ell_{\mathsf{tag}}}\bigl(P(\veck\Vert\vecn)+(\veck\Vert\vecn)\bigr),
\]
where $P$ is the Poseidon2-style permutation and $\Tr_{\ell_{\mathsf{tag}}}$ returns the first
$\ell_{\mathsf{tag}}$ lanes.
\end{definition}

\begin{assumption}[Concrete Poseidon2-style PRF assumption]\label{ass:poseidon-prf}
For the concrete parameter choices fixed in Appendix~\ref{app:prf}, the family
\[
  \PRF(k,n)=\Tr_{\ell_{\mathsf{tag}}}\bigl(P(k\Vert n)+(k\Vert n)\bigr)
\]
of Definition~\ref{def:poseidon-prf} is a secure pseudorandom function for
$k\getsunif\mathcal K_{\mathsf{key}}$.  If the signed key space is changed, this assumption is made for the
changed key distribution and the ARC relations must enforce that same key space.
\end{assumption}
